\section{Light --- Lib --- Next Pair}
%%% generated by the blueprint pipeline from TCSlib.Programs.Light.Lib.NextPair
\subsection*{Light --- Lib --- Next Pair: Introduction}

Declarations of the Lean module \texttt{TCSlib.Programs.Light.Lib.NextPair}.

\subsection*{Light --- Lib --- Next Pair: Results}

\begin{definition}[Successor pair update]
\lean{Light.nextPair}
\label{Light.nextPair}
\leanok
\uses{Light.Stmt, Light.k, Light.v}
For natural numbers $A$, $B$, and $N$, define a statement that first increments $B$ by
$1$. It then tests whether the updated value of $B$ equals $N$; if so, it sets $B$ to
$0$ and increments $A$ by $1$, and otherwise does nothing further.
\end{definition}

\begin{theorem}[Execution of a counter update]
\lean{Light.Ends.nextPair}
\label{Light.Ends.nextPair}
\leanok
\uses{Light.Cond.Holds, Light.Ends, Light.Ends.iteLast, Light.Ends.setTo, Light.Ends.setToThen, Light.Ends.skip, Light.Expr.Gives, Light.Expr.val, Light.State, Light.frame, Light.frame_setLocal, Light.nextPair, Light.setLocal, Nat.succ_div_mod_of_eq, Nat.succ_div_mod_of_ne}
\difficulty{6}
Let $A,B,N,n,t,T$ be natural numbers, let $l$ be a list of integers, let $\mu$ be an
integer-valued memory, and let $Q$ be a state predicate. Suppose $Q$ holds in the state
whose locals are the frame obtained by setting local $B$ to $(t+1)\bmod n$ and then
local $A$ to $(t+1)/n$ in $l$, and whose memory is $\mu$. If $0<n$, $t+1$ and $n$ are at
most the word limit, the initial frame has values $t/n$, $t\bmod n$, and $n$ at locals
$A,B,N$, respectively, $A\ne B$, $B\ne N$, and $14\le T$, then the statement that
increments $B$, resets it to zero and increments $A$ when the new value of $B$ equals
$N$, has an execution from the state with frame $l$ and memory $\mu$, with cost at most
$T$, ending in a state satisfying $Q$.
\end{theorem}
